%%=============================================================================
%% Methodologie
%%=============================================================================

\chapter{\IfLanguageName{dutch}{Methodologie}{Methodology}}%
\label{ch:methodologie}

%% TODO: In dit hoofstuk geef je een korte toelichting over hoe je te werk bent
%% gegaan. Verdeel je onderzoek in grote fasen, en licht in elke fase toe wat
%% de doelstelling was, welke deliverables daar uit gekomen zijn, en welke
%% onderzoeksmethoden je daarbij toegepast hebt. Verantwoord waarom je
%% op deze manier te werk gegaan bent.
%% 
%% Voorbeelden van zulke fasen zijn: literatuurstudie, opstellen van een
%% requirements-analyse, opstellen long-list (bij vergelijkende studie),
%% selectie van geschikte tools (bij vergelijkende studie, "short-list"),
%% opzetten testopstelling/PoC, uitvoeren testen en verzamelen
%% van resultaten, analyse van resultaten, ...
%%
%% !!!!! LET OP !!!!!
%%
%% Het is uitdrukkelijk NIET de bedoeling dat je het grootste deel van de corpus
%% van je bachelorproef in dit hoofstuk verwerkt! Dit hoofdstuk is eerder een
%% kort overzicht van je plan van aanpak.
%%
%% Maak voor elke fase (behalve het literatuuronderzoek) een NIEUW HOOFDSTUK aan
%% en geef het een gepaste titel.



\subsection{Fase 1: Rust aanleren}
In de eerste fase wordt de focus gelegd op het aanleren van de gekozen programmeertaal Rust. Hierbij zullen de basisconcepten van de taal geleerd worden, zoals functies,methodes,datatypes en algemene syntax. Aan de hand van verschillende kleine projecten worden deze concepten dan toegepast en aangeleerd.
\subsubsection{Week 1-2}
Tijdens deze week werd de Rust Quick Start Guide \autocite{Arbuckle2018} doorgenomen en werden er verschillende scripts uitgevoerd om diverse concepten uit te testen.





\subsection{Fase 2: Sprint 1 - opstarten}
In de eerste sprint ligt de focus op het kunnen opstarten van de virtuele machine (VM) op basis van een gemaakte testbox. De provisioner zal een VM kunnen aanmaken, configureren en starten met het virtualisatieplatform VirtualBox.


\subsubsection{Week 3}
Tijdens deze week werd het opstarten van de VM aan de hand van verschillende VBoxManage-commando's geïmplementeerd. Dit begon met het maken van een Ubuntu- en Windows 10-VM om als basisboxen te gebruiken. Via het script wordt een nieuwe VM aangemaakt met het gekozen besturingssysteem. Dit wordt gedaan door het .vdi-bestand te kopiëren van de vooraf aangemaakte VM en de Universally Unique Identifier (UUID) aan te passen, zodat er geen conflicten optreden. Vervolgens wordt deze schijf gemount en wordt de VM opgestart, nu deze klaar is voor verder gebruik.


\subsubsection{Week 4}

Tijdens deze week werd de code herschreven zodat die meer overeenkomt met Objectgeoriënteerd programmeren (OOP). Dit maakt de structuur duidelijker en vergemakkelijkt toekomstige uitbreidingen van de code.

\subsection{Fase 3: Sprint 2 - provisionen}

In de tweede stap wordt het provisioningproces verder uitgewerkt. Hierbij wordt de NAT-portforwarding aangepast zodat er via SSH verbinding kan worden gemaakt met de VM. Daarna worden de provisioningstappen uitgevoerd op het systeem.

\subsubsection{Week 5}

Tijdens deze week werd de virtuele machine verder aangepast zodat er aparte functies beschikbaar zijn voor Linux- en Windows-VM's. Hierdoor kan de uitvoering van scripts afhankelijk worden gemaakt van het type besturingssysteem.

\subsubsection{Week 6}

In deze week werd de code verder toegepast voor het verbinden met de VM via SSH en het uitvoeren van scripts op basis van het gebruikte besturingssysteem. Op die manier wordt het provisioningproces beter afgestemd op de gekozen omgeving.

\subsection{Fase 4: Sprint 3 - configuratie}
In de derde sprint wordt er ondersteuning toegevoegd voor een configuratiebestand in het TOML-formaat. De provisioner zal het configuratiebestand inlezen en verwerken, de gevraagde instellingen zullen het provisioninggedrag besturen.

\subsubsection{Week 7}

Tijdens deze week werd het lezen van het TOML-bestand opgestart. Op dat moment was de implementatie nog niet volledig uitgewerkt, maar wel al voldoende om de eerste structuur op te zetten.

\subsubsection{Week 8}

Tijdens deze week werd het inlezen van het TOML-bestand verder uitgewerkt en geïntegreerd in de werking van het programma. Daarnaast werd een init-functie toegevoegd die de juiste folderstructuur aanmaakt.



\subsection{Fase 5: Sprint 4 - functies}
In de vierde sprint zullen nieuwe functies worden toegevoegd, zoals in het aanmaken en terugzetten van snapshots. Dit zorgt er voor dat de gebruiker snel terug zal kunnen gaan naar een stabiele toestand, zoals na een mislukte provisioningrun. Ook zal tijdens deze sprint meerdere boxes aangemaakt worden.


\subsubsection{Week 9}

Tijdens deze week werd begonnen met het toepassen van snapshots. De code werd aangepast zodat het aanmaken van snapshots correct kon worden geïntegreerd in het provisioningproces. Daarnaast werden functies opgesplitst zodat ze afzonderlijk kunnen worden gebruikt voor één VM of voor meerdere VM’s. Tot slot werd de CLI-input verder verbeterd met behulp van \texttt{clap}\autocite{clapCrate}.
\subsubsection{Week 10}

Tijdens deze week werd een aparte functie toegevoegd om de status van de VM op te halen, zodat gecontroleerd kan worden of deze actief is. Daarnaast werd de bridge network adapter toegevoegd, zodat deze kan worden gebruikt wanneer de gebruiker dat wenst. Tot slot werden de snapshotfuncties verder uitgewerkt en toegepast.

\subsection{Fase 6: Sprint 5 - foutafhandeling}
In de laatste sprint ligt de focus op het bouwen van foutafhandeling en rapportering. Fouten worden gelogd met duidelijke foutboodschappen en genoteerd in een log file. Als er een fout is moet er ook op een veilige manier een fallback-actie uitgevoerd worden, dit zorgt ervoor dat de gebruiker niet met een onbruikbare omgeving blijft zitten.


\subsubsection{Week 11}
x
\subsubsection{Week 12}
x
%\newpage

\subsection{Planning}
De totale duur van deze methodologie is gepland over een periode van 12 schoolweken tijdens semester 2 van het academiejaar 2025-2026.
\begin{itemize}
    \item Week 1-2: Rust aanleren
    \item Week 3-4: Sprint 1
    \item Week 5-6: Sprint 2
    \item Week 7-8: Sprint 3
    \item Week 9-10: Sprint 4
    \item Week 11-12: Sprint 5
\end{itemize}




